\documentclass[11pt]{amsart}

\usepackage{amsmath,amssymb,amsthm,mathtools}
\usepackage[hidelinks]{hyperref}

\newtheorem{theorem}{Theorem}
\newtheorem{lemma}[theorem]{Lemma}
\theoremstyle{remark}
\newtheorem{remark}[theorem]{Remark}

\newcommand{\Q}{\mathbb Q}
\newcommand{\R}{\mathbb R}
\newcommand{\Z}{\mathbb Z}
\newcommand{\T}{\mathbb T}
\newcommand{\D}{\mathbb D}
\newcommand{\e}{\mathrm e}
\newcommand{\dd}{\,\mathrm d}

\title{A resonant four-point case of the HRT conjecture}
\author{}
\date{}

\begin{document}

\begin{abstract}
For $a\in\R\setminus\Q$ and $j\in\Z\setminus\{-1,1\}$, we prove that, for every nonzero
$g\in L^2(\R)$, the four functions
\[
 g,\qquad T_1g,\qquad M_1g,\qquad M_{a+j}T_ag
\]
are linearly independent. In particular, taking $a=\sqrt2$ and $j=0$ proves the
Heil--Speegle four-point subconjecture, and in fact proves the corresponding statement
for arbitrary nonzero $L^2$ windows. The proof combines the Zak transform and the
resonant decomposition into invariant circles with Jensen's formula and the theorem of
Iwanik--Lema\'nczyk--Rudolph on nonzero-degree cocycles over irrational rotations.
This is a result for a specific resonant family of four-point configurations; it does not
settle the HRT conjecture for arbitrary four-point sets.
\end{abstract}

\maketitle

\section{Statement}

Put $\T=\R/\Z$. For $x,\omega\in\R$, write
\[
 (T_xg)(t)=g(t-x),\qquad (M_\omega g)(t)=\e^{2\pi i\omega t}g(t).
\]
The HRT conjecture asserts that the time--frequency translates indexed by any finite set
of distinct points in $\R^2$ are linearly independent for every nonzero $g\in L^2(\R)$.
We prove the following special four-point case.

\begin{theorem}\label{thm:main}
Let $a\in\R\setminus\Q$ and let $j\in\Z\setminus\{-1,1\}$. For every nonzero
$g\in L^2(\R)$, the four functions
\[
 g,\qquad T_1g,\qquad M_1g,\qquad M_{a+j}T_ag
\]
are linearly independent.
\end{theorem}

The case $a=\sqrt2$, $j=0$ is the configuration
\[
 \Lambda_0=\{(0,0),(1,0),(0,1),(\sqrt2,\sqrt2)\}
\]
singled out by Heil and Speegle as the HRT subconjecture \cite[Conjecture~2]{HeilSpeegle}.
The argument below works directly in $L^2$; no continuity or decay assumption on $g$ is
used.

We use one external spectral theorem. If $R_\alpha t=t+\alpha$ is an irrational
rotation of $\T$ and $\xi:\T\to\T$, let
\[
 V_{\xi,\alpha}f(t)=\xi(t)f(R_\alpha t),\qquad f\in L^2(\T).
\]

\begin{theorem}[Iwanik--Lema\'nczyk--Rudolph]\label{thm:ILR}
Suppose that $\alpha$ is irrational and that $\xi:\T\to\T$ is absolutely continuous,
has nonzero topological degree, and has derivative of bounded variation. Then
$V_{\xi,\alpha}$ has Lebesgue maximal spectral type. In particular,
$V_{\xi,\alpha}$ has no eigenvectors.
\end{theorem}

This is the form of the main spectral statement in \cite{ILR}; see also
\cite[Theorem~4]{KanigowskiLemanczyk}. We will apply it only to real-analytic
circle-valued functions.

\section{Zak reduction and the invariant fibres}

Use the Zak transform normalized by
\[
 Zg(t,\omega)=\sum_{n\in\Z}g(t-n)\e^{2\pi i n\omega}.
\]
Initially this formula is interpreted on a dense subspace and then extended by unitarity
from $L^2(\R)$ onto $L^2([0,1)^2)$. Its quasiperiodic extension satisfies, almost
everywhere,
\begin{equation}\label{eq:zak-quasi}
 Zg(t+1,\omega)=\e^{2\pi i\omega}Zg(t,\omega),
 \qquad
 Zg(t,\omega+1)=Zg(t,\omega),
\end{equation}
and
\begin{equation}\label{eq:zak-covariance}
 Z(M_bT_xg)(t,\omega)=\e^{2\pi i bt}Zg(t-x,\omega-b).
\end{equation}

Assume, towards a contradiction, that
\begin{equation}\label{eq:dependence}
 A g+B T_1g+C M_1g+D M_{a+j}T_ag=0
\end{equation}
for some coefficients not all zero. If one coefficient vanishes, this is a dependence
among at most three time--frequency translates, contrary to the three-point theorem of
Heil--Ramanathan--Topiwala \cite{HRT}. Hence
\begin{equation}\label{eq:nonzero-coefficients}
 A B C D\ne0.
\end{equation}

Put $G=Zg$. By \eqref{eq:zak-quasi}, \eqref{eq:zak-covariance}, and $j\in\Z$,
\eqref{eq:dependence} becomes
\begin{equation}\label{eq:FE}
 p(t,\omega)G(t,\omega)
 =-D\e^{2\pi i(a+j)t}G(t-a,\omega-a),
\end{equation}
where
\begin{equation}\label{eq:p}
 p(t,\omega)=A+B\e^{-2\pi i\omega}+C\e^{2\pi i t}.
\end{equation}
The translation $(t,\omega)\mapsto(t-a,\omega-a)$ preserves
$\theta=\omega-t\pmod1$. For $\theta\in\T$, set
\[
 G_\theta(t)=G(t,t+\theta),
 \qquad
 P_\theta(t)=A+B\e^{-2\pi i(t+\theta)}+C\e^{2\pi i t}.
\]
The coordinate change $(t,\theta)\mapsto(t,t+\theta)$ preserves Haar measure on
$\T^2$. Consequently, for almost every $\theta$, the restriction of $G_\theta$ to
$[0,1)$ belongs to $L^2([0,1))$, and the following identities hold almost everywhere
in $t$:
\begin{align}
 G_\theta(t+1)&=\e^{2\pi i(t+\theta)}G_\theta(t),\label{eq:fibre-quasi}\\
 G_\theta(t-a)&=-D^{-1}\e^{-2\pi i(a+j)t}P_\theta(t)G_\theta(t).
 \label{eq:fibre-FE}
\end{align}
Moreover,
\begin{equation}\label{eq:fubini}
 \|g\|_2^2=\int_\T\int_0^1 |G_\theta(t)|^2\dd t\dd\theta.
\end{equation}

Call a fibre \emph{live} if $G_\theta$ is not zero as an element of $L^2([0,1))$.
Since $g\ne0$, \eqref{eq:fubini} implies that the set of live fibres has positive
measure.

\begin{lemma}[Fibre dichotomy]\label{lem:dichotomy}
For almost every $\theta$, either $G_\theta=0$ almost everywhere or
$G_\theta(t)\ne0$ for almost every $t$.
\end{lemma}

\begin{proof}
Taking absolute values in \eqref{eq:fibre-FE} gives
\begin{equation}\label{eq:modulus-cocycle}
 |G_\theta(t-a)|=|D|^{-1}|P_\theta(t)|\,|G_\theta(t)|.
\end{equation}
By \eqref{eq:fibre-quasi}, $|G_\theta|$ is $1$-periodic. The trigonometric
polynomial $P_\theta$ is not identically zero and has only finitely many zeros.
Hence $|P_\theta|>0$ almost everywhere, and \eqref{eq:modulus-cocycle} shows that
the zero set of $|G_\theta|$ is invariant modulo null sets under the irrational
rotation $t\mapsto t-a$. Ergodicity gives measure zero or one.
\end{proof}

\section{A mean condition and a counting lemma}

We first record a measurable-coboundary fact.

\begin{lemma}\label{lem:mean-zero}
Let $R$ be an ergodic measure-preserving transformation of a probability space. Suppose
$\varphi\in L^1$ and $u$ is finite and measurable almost everywhere, with
\[
 \varphi=u\circ R-u
\]
almost everywhere. Then $\int\varphi=0$.
\end{lemma}

\begin{proof}
By Birkhoff's theorem,
\[
 \frac1n\sum_{k=0}^{n-1}\varphi\circ R^k\longrightarrow\int\varphi
\]
almost everywhere. The left side equals $(u\circ R^n-u)/n$. For every
$\varepsilon>0$,
\[
 \mu\bigl(|u\circ R^n-u|>n\varepsilon\bigr)
 \le 2\mu\bigl(|u|>n\varepsilon/2\bigr)\longrightarrow0,
\]
since $u\circ R^n$ and $u$ have the same distribution. Thus the Birkhoff averages
converge to zero in measure, and hence $\int\varphi=0$.
\end{proof}

For a live fibre, Lemma~\ref{lem:dichotomy} allows us to take logarithms in
\eqref{eq:modulus-cocycle}; the values on a null set may be defined arbitrarily. Since
$\log|P_\theta|\in L^1(\T)$, Lemma~\ref{lem:mean-zero} yields
\begin{equation}\label{eq:mean-condition}
 \int_0^1\log|P_\theta(t)|\dd t=\log|D|.
\end{equation}

Write $z=\e^{2\pi i t}$ and
\begin{equation}\label{eq:Q}
 Q_\theta(z)=Cz^2+Az+B\e^{-2\pi i\theta},
 \qquad
 P_\theta(t)=z^{-1}Q_\theta(z).
\end{equation}
Let $N(\theta)$ be the number of zeros of $Q_\theta$ in the open unit disc, counted
with multiplicity, whenever $Q_\theta$ has no zero on the unit circle.

\begin{lemma}[Fixed-radius counting]\label{lem:counting}
For each $\rho>0$, there are at most two values of $\theta\in\T$ for which
$Q_\theta$ has a zero of modulus $\rho$.
\end{lemma}

\begin{proof}
Suppose that $Q_\theta(\rho\e^{i\phi})=0$. Taking absolute values gives
\[
 \rho\,|A+C\rho\e^{i\phi}|=|B|.
\]
After squaring,
\[
 \Re\!\left(C\overline A\,\e^{i\phi}\right)
 =\frac{|B|^2/\rho^2-|A|^2-|C|^2\rho^2}{2\rho}.
\]
Because $A C\ne0$, this equation has at most two solutions for $\phi$ modulo $2\pi$.
For each such zero $z=\rho\e^{i\phi}$, the equation
\[
 \e^{-2\pi i\theta}=-\frac{Cz^2+Az}{B}
\]
determines at most one $\theta\in\T$.
\end{proof}

In particular, the exceptional set
\begin{equation}\label{eq:boundary-set}
 E_\partial=\{\theta:Q_\theta\text{ has a zero on }|z|=1\}
\end{equation}
has at most two elements.

If $N(\theta)=1$, let $\zeta_{\rm out}(\theta)$ denote the zero outside the unit disc.
Jensen's formula and \eqref{eq:mean-condition} give
\begin{equation}\label{eq:jensen}
 \log|D|
 =\log|C|+\log|\zeta_{\rm out}(\theta)|.
\end{equation}
Thus every live fibre with $N(\theta)=1$ belongs to the set in
Lemma~\ref{lem:counting} with
\begin{equation}\label{eq:rho}
 \rho=\frac{|D|}{|C|}.
\end{equation}
Consequently,
\begin{equation}\label{eq:N1-finite}
 \{\theta:\theta\text{ is live and }N(\theta)=1\}
 \quad\text{has at most two elements.}
\end{equation}

\section{The phase cocycle}

Fix a live fibre $\theta\notin E_\partial$. Define
\begin{equation}\label{eq:m}
 m_\theta(t)=\e^{2\pi i\theta t}\e^{\pi i t(t-1)}.
\end{equation}
Then $|m_\theta|=1$ and
\[
 \frac{m_\theta(t+1)}{m_\theta(t)}=\e^{2\pi i(t+\theta)}.
\]
Hence
\[
 V_\theta(t)=\frac{G_\theta(t)}{m_\theta(t)}
\]
defines a measurable $1$-periodic function. A direct calculation gives
\[
 \frac{m_\theta(t)}{m_\theta(t-a)}
 =\e^{2\pi i\theta a}\e^{2\pi iat}\e^{-\pi i(a^2+a)}.
\]
Using \eqref{eq:fibre-FE}, we obtain
\begin{equation}\label{eq:periodic-cocycle}
 V_\theta(t-a)=\Lambda_\theta\e^{-2\pi ijt}P_\theta(t)V_\theta(t),
 \qquad
 \Lambda_\theta=-D^{-1}\e^{2\pi i\theta a}\e^{-\pi i(a^2+a)}.
\end{equation}

Since the fibre is live, $V_\theta\ne0$ almost everywhere. Put
\[
 W_\theta(t)=\frac{V_\theta(t)}{|V_\theta(t)|}.
\]
Because $\theta\notin E_\partial$, $P_\theta$ is zero-free on $\T$, and
\eqref{eq:periodic-cocycle} implies
\begin{equation}\label{eq:phase-cocycle}
 W_\theta(t-a)=\xi_\theta(t)W_\theta(t),
\end{equation}
where
\begin{equation}\label{eq:xi}
 \xi_\theta(t)=
 \frac{\Lambda_\theta}{|\Lambda_\theta|}
 \e^{-2\pi ijt}\frac{P_\theta(t)}{|P_\theta(t)|}.
\end{equation}
The function $\xi_\theta:\T\to\T$ is real analytic. By the argument principle,
\begin{equation}\label{eq:degree}
 \deg\xi_\theta
 =-j+\operatorname{wind}(P_\theta,0)
 =-j-1+N(\theta).
\end{equation}
Indeed, $P_\theta=z^{-1}Q_\theta$, the factor $z^{-1}$ has winding number $-1$, and
$Q_\theta(\e^{2\pi i t})$ has winding number $N(\theta)$.

\begin{lemma}\label{lem:degree-kill}
Every live fibre $\theta\notin E_\partial$ satisfies
\[
 N(\theta)-1-j=0.
\]
\end{lemma}

\begin{proof}
Suppose instead that $\deg\xi_\theta\ne0$. Let
\[
 (U_\theta f)(t)=\xi_\theta(t)f(t-a).
\]
Theorem~\ref{thm:ILR} applies to $U_\theta$, so $U_\theta$ has no eigenvectors. On the
other hand, conjugating \eqref{eq:phase-cocycle} gives
\[
 U_\theta\overline{W_\theta}
 =\xi_\theta(t)\overline{W_\theta(t-a)}
 =\overline{W_\theta(t)}.
\]
Thus $\overline{W_\theta}$ is a nonzero eigenvector with eigenvalue $1$, a contradiction.
\end{proof}

\section{Conclusion}

\begin{proof}[Proof of Theorem~\ref{thm:main}]
Suppose first that $|j|\ge2$. Since $N(\theta)\in\{0,1,2\}$,
\[
 N(\theta)-1-j\ne0
\]
for every $\theta\notin E_\partial$. Lemma~\ref{lem:degree-kill} therefore shows that
every live fibre belongs to the finite set $E_\partial$, up to a null set. This
contradicts the fact that the live fibres have positive measure.

It remains to treat $j=0$. Lemma~\ref{lem:degree-kill} excludes every live fibre
outside $E_\partial$ with $N(\theta)=0$ or $2$. The live fibres with $N(\theta)=1$
form a set of at most two points by \eqref{eq:N1-finite}. Thus the live-fibre set is
again finite up to a null set, contradicting \eqref{eq:fubini}.

Hence no nontrivial dependence \eqref{eq:dependence} exists.
\end{proof}

\begin{remark}
For $j=1$, the phase degree can vanish when both zeros of $Q_\theta$ lie in the unit
disc; for $j=-1$, it can vanish when both lie outside. The argument above therefore
stops exactly at $j=\pm1$. No claim about those two families is made here.
\end{remark}

\begin{remark}[Scope]
Theorem~\ref{thm:main} proves one resonant family of four-point configurations. It is
not a proof of the HRT conjecture for every set of four points. In particular, the
condition that the fourth point have the form $(a,a+j)$ is essential to the invariant
circle decomposition used above.
\end{remark}

\begin{remark}[Relation to prior work]
The Zak reduction and the decomposition of the resonant mixed-integer case into
invariant cosets are part of Oussa's trichotomy program \cite{Oussa}. The present proof
uses, in addition, the fixed-radius consequence of Jensen's formula and the spectral
obstruction supplied by Theorem~\ref{thm:ILR}. Oussa's current arXiv version derives
necessary dynamical, cohomological, and arithmetic constraints in the infinite proper
orbit case, rather than claiming a proof of that case.
\end{remark}

\begin{thebibliography}{99}

\bibitem{HRT}
C.~Heil, J.~Ramanathan, and P.~Topiwala,
\emph{Linear independence of time--frequency translates},
Proc. Amer. Math. Soc. \textbf{124} (1996), 2787--2795.

\bibitem{HeilSpeegle}
C.~Heil and D.~Speegle,
\emph{The HRT conjecture and the zero divisor conjecture for the Heisenberg group},
in \emph{Excursions in Harmonic Analysis}, Vol.~3,
Birkh\"auser, 2015, 159--176.

\bibitem{ILR}
A.~Iwanik, M.~Lema\'nczyk, and D.~Rudolph,
\emph{Absolutely continuous cocycles over irrational rotations},
Israel J. Math. \textbf{83} (1993), 73--95.

\bibitem{KanigowskiLemanczyk}
A.~Kanigowski and M.~Lema\'nczyk,
\emph{Spectral theory of dynamical systems},
arXiv:2006.11616 (2020).

\bibitem{Oussa}
V.~Oussa,
\emph{Trichotomy for the HRT conjecture for mixed integer configuration},
arXiv:2508.04613v2 (2026).

\end{thebibliography}

\end{document}
